% =============================================================================
\section{Additional Binding Tools}
\label{sec:tools}
% =============================================================================
In this section we list several tools that facilitate binding generation, in addition to the tools listed in Section~\ref{sec:comp}.

% ctypes
\pyLstinline{ctypes} is a foreign function module for Python that provides C compatible data types.\footnote{See \url{https://docs.python.org/3/library/ctypes.html\#module-ctypes}.} It supports the loading of shared libraries and marshaling data between Python and C.\footnote{Marshaling is the process of transforming the memory representation of an object into a data format suitable for storage or transmission; see \url{https://en.wikipedia.org/wiki/Marshalling\_(computer\_science)}.} The \pyLstinline{ctypes} module is part of the Python standard library; thus, it can be used to wrap shared libraries in pure Python. Using the \pyLstinline{ctypes} module requires manual programming labour and detailed knowledge of the \pyLstinline{ctypes} interface. However, it does not require writing C code (in addition to existing C code that needs to be wrapped).

% CFFI
\pyLstinline{cffi} is another Python module that generates Python bindings (it stands for C Foreign Function Interface). It interacts with almost any C code from Python, based on C-like declarations that can often be copy-pasted from header files. It is compatible with both PyPy and CPython; Moreover, in the PyPy (see~\ref{ssec:intro:bind})
reference implementation it is a first class citizen (and provided by defaut), thus, exhibits outstanding performance. Compared to \pyLstinline{ctypes}, the \pyLstinline{cffi} module takes a more automated approach to generate Python bindings and it minimizes the extra bits of API that needs to be mastered; thus, it scales better to larger projects than \pyLstinline{ctypes}. Like \pyLstinline{ctypes}, using \pyLstinline{cffi} directly interfaced with C libraries only. Wrapping C++ libraries with any one of these modules requires the provision of another layer of C wrapper around the C++ code.

% cppyy
\pyLstinline{cppyy} is yet another external module for generating bindings.\footnote{See \url{https://cppyy.readthedocs.io}.} It is an automatic, run-time, C++ Python bindings generator, for calling C++ from Python and Python from C++.\footnote{See \url{https://cppyy.readthedocs.io}.} Run-time generation enables detailed specialization for higher performance, lazy loading for reduced memory use in large scale projects,
Python-side cross-inheritance and callbacks for working with C++ frameworks, run-time template instantiation, automatic object downcasting, exception mapping, and interactive exploration of C++ libraries. \pyLstinline{cppyy} delivers this without any language
extensions, intermediate languages, or the need for boiler-plate hand-written code. For design and performance, see~\cite{ld-hppcn-16}, albeit that the \pyLstinline{cppyy} performance has been vastly improved since.  Like \pyLstinline{cffi}, \pyLstinline{cppyy} was designed to ease binding development for Python programmer,
minimizing the need for extra C++ code; it is also compatible with both CPython and PyPy.

% Cython
A completely different approach is adopted by Cython~\cite{bbcds-cbbw-11}, which refers to a programming language and to an optimising static compiler for the Cython language.\footnote{See \url{https://cython.org/}.}  The Cython
language is a superset of the Python language. It extends Python with explicit type declarations of native C/C++ types.  Cython is a compiled language used to generate CPython extension modules.  The Cython compiler generates C/C++ code from Cython code, which in turn compiles with any C/C++ compiler, and automatically wrapped in interface code, producing extension modules that can be loaded and used by regular Python code using the \pyLstinline{import} statement, but with significantly less computational overhead at run time. Cython also facilitates wrapping independent C or C++ code into
python-importable modules.  With Cython it is possible to call back and forth between Python code, Cython code, and native library code originally developed in Fortrant, C, or C++.

% SIP
SIP is a collection of tools that makes it easy to create Python bindings for C and C++ libraries.\footnote{See \url{https://pypi.org/project/sip/}.}  It was originally developed in 1998 to create PyQt, the Python bindings for the Qt toolkit, but can be used to create bindings for any C or C++ library. SIP comprises a set of build tools and a module called sip. The build tools process a set of specification files and generates C or C++ code, which is then compiled to create the bindings extension module. Several extension modules may be installed in the same Python package. Extension modules can be built so that they are are independent of the version of Python being used.
% In other words a wheel created from them can be installed with any
% version of Python starting with v3.5.

% PyBindGen
